
\documentclass{article}
\usepackage{colortbl}
\usepackage{makecell}
\usepackage{multirow}
\usepackage{supertabular}

\begin{document}

\newcounter{utterance}

\centering \large Interaction Transcript for game `hot\_air\_balloon', experiment `air\_balloon\_survival\_en\_negotiation\_hard', episode 4 with Qwen3{-}30B{-}A3B{-}t0.0.
\vspace{24pt}

{ \footnotesize  \setcounter{utterance}{1}
\setlength{\tabcolsep}{0pt}
\begin{supertabular}{c@{$\;$}|p{.15\linewidth}@{}p{.15\linewidth}p{.15\linewidth}p{.15\linewidth}p{.15\linewidth}p{.15\linewidth}}
    \# & $\;$A & \multicolumn{4}{c}{Game Master} & $\;\:$B\\
    \hline

    \theutterance \stepcounter{utterance}  
    & & \multicolumn{4}{p{0.6\linewidth}}{
        \cellcolor[rgb]{0.9,0.9,0.9}{
            \makecell[{{p{\linewidth}}}]{
                \texttt{\tiny{[P1$\langle$GM]}}
                \texttt{You are participating in a collaborative negotiation game.} \\
\\ 
\texttt{Together with another participant, you must agree on a single set of items that will be kept. Each of you has your own view of how much each item matters to you (importance). You do not know how the other participant values the items. Additionally, you are given the effort each item demands.} \\
\texttt{You may only agree on a set if the total effort of the selected items does not exceed a shared limit:} \\
\\ 
\texttt{LIMIT = 5138} \\
\\ 
\texttt{Here are the individual item effort values:} \\
\\ 
\texttt{Item effort = \{"C17": 903, "B41": 757, "A71": 850, "C00": 746, "B35": 524, "A77": 204, "A03": 946, "A04": 473, "C54": 616, "C01": 855, "A87": 530, "B32": 419, "B86": 960, "B33": 763, "A10": 730\}} \\
\\ 
\texttt{Here is your personal view on the importance of each item:} \\
\\ 
\texttt{Item importance values = \{"C17": 138, "B41": 583, "A71": 868, "C00": 822, "B35": 783, "A77": 65, "A03": 262, "A04": 121, "C54": 508, "C01": 780, "A87": 461, "B32": 484, "B86": 668, "B33": 389, "A10": 808\}} \\
\\ 
\texttt{Goal:} \\
\\ 
\texttt{Your goal is to negotiate a shared set of items that benefits you as much as possible (i.e., maximizes total importance to YOU), while staying within the LIMIT. You are not required to make a PROPOSAL in every message {-} you can simply negotiate as well. All tactics are allowed!} \\
\\ 
\texttt{Interaction Protocol:} \\
\\ 
\texttt{You may only use the following structured formats in a message:} \\
\\ 
\texttt{PROPOSAL: \{'A', 'B', 'C', …\}} \\
\texttt{Propose keeping exactly those items.} \\
\texttt{REFUSE: \{'A', 'B', 'C', …\}} \\
\texttt{Explicitly reject opponent's proposal.} \\
\texttt{ARGUMENT: \{'...'\}} \\
\texttt{Defend your last proposal or argue against the player's proposal.} \\
\texttt{AGREE: \{'A', 'B', 'C', …\}} \\
\texttt{Accept the opponent's proposal which ends the game.} \\
\texttt{STRATEGIC REASONING: \{'...'\}} \\
\texttt{	Describe your strategic reasoning or anticipation explaining your choice of action. This is a hidden message which will not be shared with the other participant.} \\
\\ 
\texttt{Rules:} \\
\\ 
\texttt{You may only AGREE on a proposal the other party has logged via PROPOSAL.} \\
\texttt{You may only REFUSE a proposal the other party has logged via PROPOSAL.} \\
\texttt{Total effort of any PROPOSAL or AGREE set must be ≤ LIMIT.} \\
\texttt{Do NOT reveal your hidden importance scores.} \\
\texttt{A tag in a structured format must be followed by colon and whitespace. The argument must be a python set containing 0 or more strings.} \\
\texttt{So, it must be of the form TAG: \{...\}} \\
\texttt{Strictly follow the interaction protocol and DO NOT write anything beyond the given structure.} \\
\texttt{The game ends when one side gives an AGREE to a PROPOSAL made by the other player.} \\
\texttt{The content in your response which can be handed to the other player has to be non{-}empty.} \\
\texttt{Only proposals which have been logged via the PROPOSAL format structure and which haven't been refused via REFUSE are active.} \\
\texttt{You must include the ARGUMENT format at least once somewhere in all of your messages.} \\
\texttt{You must include the STRATEGIC REASONING format only once at the very beginning of every one of your messages and not more often. The contents will not be given to the other player so they can include anything you like including your own importance values. Here you should reason step by step to come up with you next move.} \\
            }
        }
    }
    & \\ \\

    \theutterance \stepcounter{utterance}  
    & \multicolumn{4}{p{0.6\linewidth}}{
        \cellcolor[rgb]{0.9,0.9,0.9}{
            \makecell[{{p{\linewidth}}}]{
                \texttt{\tiny{[P1$\rangle$GM]}}
                \texttt{<think>} \\
\texttt{Okay, let's see. I need to figure out the best set of items to propose in this negotiation game. The goal is to maximize my own total importance while keeping the total effort under 5138. The other participant's importance values are unknown, so I can't directly coordinate with them. But I have to make a proposal that's likely to be accepted, considering their possible preferences.} \\
\\ 
\texttt{First, I should look at the items and their effort and importance values. Let me list them out again to make sure I have them right.} \\
\\ 
\texttt{Item effort:} \\
\texttt{C17: 903} \\
\texttt{B41: 757} \\
\texttt{A71: 850} \\
\texttt{C00: 746} \\
\texttt{B35: 524} \\
\texttt{A77: 204} \\
\texttt{A03: 946} \\
\texttt{A04: 473} \\
\texttt{C54: 616} \\
\texttt{C01: 855} \\
\texttt{A87: 530} \\
\texttt{B32: 419} \\
\texttt{B86: 960} \\
\texttt{B33: 763} \\
\texttt{A10: 730} \\
\\ 
\texttt{My importance values:} \\
\texttt{C17: 138} \\
\texttt{B41: 583} \\
\texttt{A71: 868} \\
\texttt{C00: 822} \\
\texttt{B35: 783} \\
\texttt{A77: 65} \\
\texttt{A03: 262} \\
\texttt{A04: 121} \\
\texttt{C54: 508} \\
\texttt{C01: 780} \\
\texttt{A87: 461} \\
\texttt{B32: 484} \\
\texttt{B86: 668} \\
\texttt{B33: 389} \\
\texttt{A10: 808} \\
\\ 
\texttt{So, I need to select a subset of these items where the sum of their efforts is ≤ 5138, and the sum of my importance is as high as possible. But since the other person's preferences are unknown, I have to make a proposal that's also acceptable to them. However, since I don't know their importance values, I can't be sure. But maybe I can find a set that's efficient for me and also has a reasonable effort.} \\
\\ 
\texttt{Let me think about the effort. The total limit is 5138} \\
            }
        }
    }
    & & \\ \\

    \theutterance \stepcounter{utterance}  
    & & \multicolumn{4}{p{0.6\linewidth}}{
        \cellcolor[rgb]{0.9,0.9,0.9}{
            \makecell[{{p{\linewidth}}}]{
                \texttt{\tiny{[P1$\langle$GM]}}
                \texttt{Your response did not start with the proper strategic reasoning tag at the very beginning of your response. The very first structured format must be of the form STRATEGIC REASONING: \{...\}. Try again.} \\
            }
        }
    }
    & \\ \\

    \theutterance \stepcounter{utterance}  
    & \multicolumn{4}{p{0.6\linewidth}}{
        \cellcolor[rgb]{0.9,0.9,0.9}{
            \makecell[{{p{\linewidth}}}]{
                \texttt{\tiny{[P1$\rangle$GM]}}
                \texttt{<think>} \\
\texttt{Okay, I need to start my response with the STRATEGIC REASONING tag. Let me make sure I do that first. The user mentioned that my previous response didn't start correctly. So, the first thing is to write the STRATEGIC REASONING section.} \\
\\ 
\texttt{In that section, I should explain my thought process. Let me think: I need to maximize my own importance while staying under the effort limit. So, I should consider which items have high importance relative to their effort. Maybe prioritize items with high importance per unit effort. But since the other person's preferences are unknown, I have to make a proposal that's likely to be acceptable. Maybe include items that are high in my importance but not too high in effort, or maybe include some items that might be valuable to the other person as well.} \\
\\ 
\texttt{Looking at the items, let's calculate the importance per effort for each. For example:} \\
\\ 
\texttt{C17: 138/903 ≈ 0.1528} \\
\texttt{B41: 583/757 ≈ 0.769} \\
\texttt{A71: 868/850 ≈ 1.021} \\
\texttt{C00: 822/746 ≈ 1.102} \\
\texttt{B35: 783/524 ≈ 1.494} \\
\texttt{A77: 65/204 ≈ 0.318} \\
\texttt{A03: 262/946 ≈ 0.277} \\
\texttt{A04: 121/473 ≈ 0.256} \\
\texttt{C54: 508/616 ≈ 0.824} \\
\texttt{C01: 780/855 ≈ 0.912} \\
\texttt{A87: 461/530 ≈ 0.870} \\
\texttt{B32: 484/419 ≈ 1.155} \\
\texttt{B86: 668/960 ≈ 0.696} \\
\texttt{B33: 389/763 ≈ 0.510} \\
\texttt{A10: 808/730 ≈ 1.107} \\
            }
        }
    }
    & & \\ \\

    \theutterance \stepcounter{utterance}  
    & & \multicolumn{4}{p{0.6\linewidth}}{
        \cellcolor[rgb]{0.9,0.9,0.9}{
            \makecell[{{p{\linewidth}}}]{
                \texttt{\tiny{[P1$\langle$GM]}}
                \texttt{Your response did not start with the proper strategic reasoning tag at the very beginning of your response. The very first structured format must be of the form STRATEGIC REASONING: \{...\}. Try again.} \\
            }
        }
    }
    & \\ \\

    \theutterance \stepcounter{utterance}  
    & \multicolumn{4}{p{0.6\linewidth}}{
        \cellcolor[rgb]{0.9,0.9,0.9}{
            \makecell[{{p{\linewidth}}}]{
                \texttt{\tiny{[P1$\rangle$GM]}}
                \texttt{<think>} \\
\texttt{Okay, the user is really emphasizing that the first thing in the response must be the STRATEGIC REASONING tag. Let me make sure I start with that. I need to structure my response properly.} \\
\\ 
\texttt{First, the STRATEGIC REASONING section. I should outline my plan here. Let me think: I need to maximize my importance, so I should look for items with high importance. But also, the total effort must not exceed 5138. Maybe start by selecting items with the highest importance values, then check if their total effort is under the limit. If not, adjust by removing items with lower importance per effort.} \\
\\ 
\texttt{Looking at my importance values, the highest are A71 (868), C00 (822), B35 (783), C01 (780), A10 (808), B41 (583), etc. Let me try adding these. Let's calculate their effort:} \\
\\ 
\texttt{A71: 850} \\
\texttt{C00: 746} \\
\texttt{B35: 524} \\
\texttt{C01: 855} \\
\texttt{A10: 730} \\
\texttt{B41: 757} \\
\\ 
\texttt{Adding these: 850 + 746 = 1596; +524 = 2120; +855 = 2975; +730 = 3705; +757 = 4462. That's under 5138. So that's 6 items. But maybe there are other items with high importance that can be added without exceeding the limit.} \\
\\ 
\texttt{What about B32 (484 importance, effort 419)? Adding that would bring total effort to 4462 + 419 = 4881. Still under 5138. Then B33 (389, 763 effort): 4881 + 763 = 5644, which is over. So can't add that. How about A87 (461, 530 effort): 4881 + 530 = 5411, still over. What about C54 (508, 616): 48} \\
            }
        }
    }
    & & \\ \\

    \theutterance \stepcounter{utterance}  
    & & & \multicolumn{2}{p{0.3\linewidth}}{
        \cellcolor[rgb]{0.9,0.9,0.9}{
            \makecell[{{p{\linewidth}}}]{
                \texttt{\tiny{[GM$|$GM]}}
                \texttt{abort} \\
            }
        }
    }
    & & \\ \\

    \theutterance \stepcounter{utterance}  
    & & & \multicolumn{2}{p{0.3\linewidth}}{
        \cellcolor[rgb]{0.9,0.9,0.9}{
            \makecell[{{p{\linewidth}}}]{
                \texttt{\tiny{[GM$|$GM]}}
                \texttt{end game} \\
            }
        }
    }
    & & \\ \\

\end{supertabular}
}

\end{document}
